%% RESUMO EM LINGUA ESTRANGEIRA (elemento obrigatorio -- NBR 6028)
%% Traducao fiel do resumo em portugues. As keywords sao definidas em
%% config/dados.tex.

\begin{abstract}
Retrieval-Augmented Generation (RAG) systems depend on retrieval quality, and
over Brazilian legal-administrative corpora, pure vector search fails at
retrieving exact terms and at answering questions that connect distinct
documents. This work empirically compares retrieval strategies --- dense RAG,
dense-sparse RAG and incremental GraphRAG, implemented on PostgreSQL with
pgvector and Apache AGE as a single store --- and generation strategies ---
single-pass and multi-agent verified generation, following the
Researcher-Auditor-Adjudicator pattern --- over a corpus of Brazilian
normative documents: federal laws, resolutions of the National Council of
Justice and normative instructions of the Federal Court of Accounts. The
evaluation uses a golden dataset of intra- and cross-document questions, the
RAGAS metrics (Faithfulness, Context Precision, Context Recall and Answer
Relevancy) and the paired Wilcoxon signed-rank test with Bonferroni
correction. In the confirmatory phase, incremental GraphRAG outperforms dense
RAG but not dense-sparse RAG, and multi-agent generation improves neither
Faithfulness nor the citation verification rate over single-pass generation:
hypotheses H1 and H2 are not confirmed. An exploratory phase extended the
comparison to fourteen retrieval conditions and showed that removing the
reranking step from dense-sparse RAG raises Context Precision with no loss on
the other metrics, and that pure sparse search (BM25) achieves the highest
Context Precision among the formal conditions. The most precise configuration
found combines dense-sparse RAG without reranking with single-pass
generation, which is simpler and cheaper than the one the hypotheses
predicted. Pure GraphRAG, which hands the generator only synthetic entity
descriptions, performed below the other conditions, and the explanation best
supported by the data is a mismatch between that format and the Faithfulness
metric. Among the study's limitations are the use of a single LLM judge and a
retrieval metric based on chunk identifiers rather than character overlap.
\end{abstract}
